\section{Occupation, Mellin Dilation, and Residue Hilbert Spaces}
\label{v4-arith:mo:chapter}

\subsection{Occupation states and two Mellin actions}

An energy proportional to $\log m$ can be represented by a point
mass at $m^{-1}$. The Mellin transform of that mass is a Dirichlet
monomial. This gives an exact map from occupation states to
distributions. Its behavior at a zero of an entire function can
therefore be calculated before choosing a Hilbert norm on the
residue values.

All vector spaces are over $\mathbb C$. Inner products are linear
in the first variable. The coordinate $x>0$ is dimensionless,
and $t=\log x$ uses the real logarithm. Put
$\mathscr D=C_c^\infty(\mathbb R_{>0})$. Its elements are smooth
functions supported in a compact subset of the open half-line.
For $f\in\mathscr D$, define
\[
 F(s)=\mathcal Mf(s)=\int_0^\infty f(x)x^{s-1}\,dx.
\]
Compact support away from zero makes $F$ entire and permits every
spectral derivative under the integral. Define two operators on
this common invariant domain by
\[
 Qf=-(\log x)f,\qquad Df=x f'(x).
\]
Differentiation under the integral and integration by parts give
\begin{equation}
 \mathcal M(Qf)=-F',\qquad \mathcal M(Df)=-sF.
 \label{v4-arith:mo:two-actions}
\end{equation}
The boundary term $[f(x)x^s]_0^\infty$ is zero. Thus $Q$ and $D$
already have different actions on entire Mellin transforms.

\subsubsection{An exact occupation embedding}

Let $\mathscr E=c_{00}(\mathbb N_{\geq1})$ be the space of finite
sequences, with basis $v_m$ and inner product
$\langle a,b\rangle=\sum_m a_m\overline{b_m}$. Define
\[
 Nv_m=(\log m)v_m.
\]
Its Hilbert completion is $\ell^2(\mathbb N_{\geq1})$. The maximal
self-adjoint extension of $N$ has domain
\[
 \operatorname{Dom}N=
 \left\{a:\sum_m(\log m)^2|a_m|^2<\infty\right\}.
\]
Indeed finite truncations converge in the graph norm, and the
adjoint equation on the coordinate vectors imposes exactly this
domain. The coefficients $(\log m\pm i)^{-1}$ give everywhere
bounded inverses for $N\pm i$.

A distribution acts complex-linearly on smooth test functions.
The Dirac distribution $\delta_a$ satisfies
$\langle\delta_a,\phi\rangle=\phi(a)$; this convention uses the
additive measure $dx$. A compactly supported distribution $u$
acts on $x^{s-1}$ after insertion of a smooth cutoff equal to one
near its support. The result is independent of the cutoff. This
defines its entire Mellin transform.

\begin{proposition}[occupation distributions]
\label{v4-arith:mo:occupation}
The map
\begin{equation}
 T:\mathscr E\longrightarrow\mathscr E'(\mathbb R_{>0}),
 \qquad Tv_m=m^{-1}\delta_{m^{-1}}
 \label{v4-arith:mo:embedding}
\end{equation}
is injective. Here $\mathscr E'$ denotes compactly supported
distributions. Multiplication of distributions by the smooth
function $-\log x$ satisfies
\[
 QT=TN,\qquad
 \mathcal M(Ta)(s)=\sum_m a_m m^{-s}.
\]
\end{proposition}
\begin{proof}
A bump supported near $m^{-1}$, excluding the finitely many other
support points, recovers $a_m/m$ from $Ta$. This proves injectivity.
On that point mass, $-\log x$ takes the value $\log m$.
Finally,
$m^{-1}(m^{-1})^{s-1}=m^{-s}$, proving the Mellin formula.
\end{proof}

The factor $m^{-1}$ is fixed by the additive-distribution convention.
If point masses were instead introduced relative to $dx/x$, their
conversion to additive distributions would contain this same factor.
The weighted occupation lattice with coordinate vectors $d_m$ of
squared norm $m^{-1}$ maps isometrically to the occupation space by
$d_m\mapsto m^{-1/2}v_m$. Its resulting distribution is
$m^{-3/2}\delta_{m^{-1}}$, with Mellin transform $m^{-s-1/2}$.

\begin{proposition}[extension and the analytic half-plane]
\label{v4-arith:mo:dirichlet-domain}
The same locally finite sum defines an injective map
$T:\ell^2\to\mathscr D'(\mathbb R_{>0})$. It is continuous for
the strong distribution topology, and $QT=TN$ on
$\operatorname{Dom}N$. For $a\in\ell^2$, the Dirichlet series
\[
 F_a(s)=\sum_m a_m m^{-s}
\]
and all its spectral derivatives converge absolutely and locally
uniformly on $\operatorname{Re}s>1/2$. No continuation across the
line $\operatorname{Re}s=1/2$ is part of this construction.
\end{proposition}
\begin{proof}
A compact subset of $\mathbb R_{>0}$ meets only finitely many points
$m^{-1}$. Thus the sum acts on every test function. Bounded sets
of tests have a common compact support and uniform bounds on all
derivatives. Only a fixed finite set of coefficients occurs on such
a set, so Cauchy--Schwarz proves strong continuity. Separate bumps
still recover every coefficient, and multiplication proves the
operator identity.

For $\sigma>1/2$ and integer $k\geq0$,
\[
 \sum_m |a_m|(\log m)^k m^{-\sigma}
 \leq\|a\|_2
 \left(\sum_m(\log m)^{2k}m^{-2\sigma}\right)^{1/2}<\infty.
\]
Integral comparison proves convergence, uniformly when $\sigma$
stays above $1/2+\epsilon$. Termwise differentiation is therefore
valid. The estimate supplies no holomorphic neighborhood of a
point on the boundary line.
\end{proof}

For a distribution $u$, set $\langle u',\phi\rangle=-\langle
u,\phi'\rangle$ and $Du=xu'$. The same calculation gives
$\mathcal M(Du)=-s\mathcal Mu$ for compactly supported $u$.
In particular,
\[
 D(Tv_m)=m^{-1}\bigl(m^{-1}\delta'_{m^{-1}}-\delta_{m^{-1}}\bigr).
\]
Dilation does not preserve the span of the point masses in
\eqref{v4-arith:mo:embedding}. The identity $QT=TN$ cannot be replaced by
an identity involving $D$.

\subsubsection{Residues of occupation distributions}

Let $\Xi$ be a nonzero entire function. Fix a nonempty finite set
$Z$ of distinct zeros. For $\rho\in Z$, let $r_\rho\geq1$ be its
multiplicity, and write
\[
 \Xi(s)^{-1}=
 \sum_{\ell=-r_\rho}^{\infty}b_{\rho,\ell}(s-\rho)^\ell,
 \qquad b_{\rho,-r_\rho}\ne0.
\]
For an entire $F$, define its residue coordinates by
\begin{equation}
 R_{\rho,k}F=\operatorname*{Res}_{s=\rho}
 \frac{(s-\rho)^kF(s)}{\Xi(s)},
 \quad 0\leq k<r_\rho.
 \label{v4-arith:mo:jets}
\end{equation}
The residue is the coefficient of $(s-\rho)^{-1}$. These coordinates
record the full principal part of $F/\Xi$. Direct multiplication
of the Taylor and Laurent series gives
\begin{equation}
 R_{\rho,k}F=
 \sum_{j=0}^{r_\rho-k-1}
 b_{\rho,-k-j-1}\frac{F^{(j)}(\rho)}{j!}.
 \label{v4-arith:mo:jet-formula}
\end{equation}
The transformation from the Taylor jet of order $r_\rho-1$ to
these coordinates is invertible, since its successive leading
coefficients are $b_{\rho,-r_\rho}$.

\begin{proposition}[the occupation residue defect]
\label{v4-arith:mo:occupation-defect}
Define $A_0a=(R_{\rho,0}F_a)_{\rho\in Z}$ on $\mathscr E$.
Let $K_\lambda$ multiply the $\rho$ coordinate by a prescribed
number $\lambda_\rho$. Then $A_0N=K_\lambda A_0$ fails on
$\mathscr E$ for every choice of the numbers $\lambda_\rho$.
At a simple zero,
\begin{equation}
 (A_0v_m)_\rho=\frac{m^{-\rho}}{\Xi'(\rho)},\qquad
 ((A_0N-K_\lambda A_0)v_m)_\rho
 =\frac{(\log m-\lambda_\rho)m^{-\rho}}{\Xi'(\rho)}.
 \label{v4-arith:mo:simple-occupation}
\end{equation}
\end{proposition}
\begin{proof}
Formula \eqref{v4-arith:mo:jet-formula} yields
\[
 (A_0v_m)_\rho=m^{-\rho}p_\rho(\log m),\qquad
 p_\rho(u)=\sum_{j=0}^{r_\rho-1}
 b_{\rho,-j-1}\frac{(-u)^j}{j!}.
\]
The polynomial has degree $r_\rho-1$ and nonzero leading
coefficient. Among $m=1,\ldots,r_\rho+1$, at least two values
$p_\rho(\log m)$ are nonzero. The proposed identity would make
$\lambda_\rho$ equal to two distinct logarithms. At a simple zero,
$m=1$ and $m=2$ already decide the question.
\end{proof}

For the local double-zero model $\Xi(s)=(s-\rho)^2$,
\[
 R_{\rho,0}(m^{-s})=-(\log m)m^{-\rho},\qquad
 R_{\rho,1}(m^{-s})=m^{-\rho}.
\]
The scalar residue kills $v_1$, but $v_2$ and $v_3$ still contradict
a fixed diagonal energy. This local entire function is a model
of the multiplicity calculation, not an identification with an
arithmetic entire function.

For general $a\in\mathscr E$, equation \eqref{v4-arith:mo:two-actions} gives
\[
 (A_0Na)_\rho=\operatorname*{Res}_{s=\rho}
                      \frac{-F_a'(s)}{\Xi(s)}.
\]
Thus allowing point distributions repairs the smooth-support
obstruction to $T$ itself. It leaves the derivative in the residue
composition. The target $\lambda_\rho=\operatorname{Im}\rho$ has
exactly the defect in \eqref{v4-arith:mo:simple-occupation}.

\subsubsection{A one-frequency obstruction before completion}

For real $c$, put
\[
 \mathscr H_c=L^2(\mathbb R_{>0},x^{2c-1}\,dx),\qquad
 (U_cf)(t)=e^{ct}f(e^t).
\]
Change of variables gives $\|U_cf\|_{L^2(dt)}=\|f\|_{\mathscr H_c}$.
Fix a real frequency $\gamma_0$ and choose
$\psi\in C_c^\infty(\mathbb R)$ with $\int\psi(t)\,dt=1$.
For $L>0$, define
\[
 g_L(t)=L^{-1}\psi(t/L)e^{-i\gamma_0t},\qquad f_L=U_c^{-1}g_L.
\]
Then $f_L\in\mathscr D$ and
\begin{equation}
 \|f_L\|_{\mathscr H_c}^2=L^{-1}\|\psi\|_2^2\longrightarrow0,
 \qquad \mathcal Mf_L(c+i\gamma_0)=1.
 \label{v4-arith:mo:continuous-sequence}
\end{equation}
For a simple zero $\rho=c+i\gamma_0$, its residue coordinate
therefore remains $1/\Xi'(\rho)$ while the source norm tends to zero.
A linear operator is \emph{closable} if every sequence of domain
vectors tending to zero whose images converge has image limit zero.
This necessary and sufficient condition says that the closure of
its graph is still the graph of a function. Equation
\eqref{v4-arith:mo:continuous-sequence} proves that this one-coordinate
evaluation is not closable on $\mathscr H_c$.

There is an equally explicit occupation sequence. Let
$H_M=\sum_{m=1}^M m^{-1}$ and fix $\rho=1/2+i\gamma_0$. Set
\[
 a_m^{(M)}=
 \begin{cases}H_M^{-1}m^{-1/2+i\gamma_0},&m\leq M,\\0,&m>M.
 \end{cases}
\]
Then
\begin{equation}
 \|a^{(M)}\|_{\ell^2}^2=H_M^{-1}\longrightarrow0,
 \qquad F_{a^{(M)}}(\rho)=1.
 \label{v4-arith:mo:atomic-sequence}
\end{equation}
Divergence of $H_M$ follows from comparison with $\int_1^{M+1}dx/x$.
Consequently one critical-line evaluation, and hence its nonzero
simple-residue multiple as a scalar operator, is not closable on
occupation $\ell^2$. This statement concerns that scalar operator.
It does not infer nonclosability of an arbitrary vector-valued
map merely from one unbounded coordinate.

\begin{proposition}[simultaneous critical-line occupation evaluations]
\label{v4-arith:mo:finite-occupation-graph}
Let $\rho_j=1/2+i\gamma_j$, $1\leq j\leq n$, have distinct
real ordinates. The map
$a\mapsto(F_a(\rho_j))_{j=1}^n$ on $\mathscr E$ has dense graph
in $\ell^2\oplus\mathbb C^n$. The same holds after multiplication
of its coordinates by any fixed nonzero constants, including
simple-residue factors.
\end{proposition}
\begin{proof}
Let $S_M$ be the matrix with entries $(S_M)_{jm}=m^{-\rho_j}$,
$1\leq m\leq M$. Its Gram matrix $G_M=S_MS_M^*$ has diagonal
entries $H_M$. For $j\ne k$, its entry is
$\sum_{m\leq M}m^{-1+i(\gamma_k-\gamma_j)}$, which is bounded
uniformly in $M$. Indeed its integral is bounded by
$2/|\gamma_k-\gamma_j|$, and the difference between the sum and
integral is bounded by the integral of the absolute derivative,
which is a constant times $\int_1^\infty x^{-2}\,dx$.
Thus $G_M=H_MI+O(1)$ in the fixed matrix dimension. For large $M$
it is invertible and $\|G_M^{-1}\|\to0$.
Given $y\in\mathbb C^n$, extend
$a^{(M)}=S_M^*G_M^{-1}y$ by zero outside $m\leq M$. Then
\[
 S_Ma^{(M)}=y,\qquad
 \|a^{(M)}\|_2^2=y^*G_M^{-1}y\longrightarrow0.
\]
The graph closure contains every $(0,y)$. Subtracting such
vectors from the graph over the dense finite-sequence core proves
density in the entire product.
\end{proof}

\begin{proposition}[finite occupation jets]
\label{v4-arith:mo:occupation-jet-graph}
For finitely many distinct zeros on $\operatorname{Re}s=1/2$,
the full residue map on $\mathscr E$ has dense graph in
$\ell^2\oplus\mathscr J_Z$. Its range is all of $\mathscr J_Z$.
For any nonempty finite zero set, its kernel on $\mathscr E$
is not invariant under $N$.
\end{proposition}
\begin{proof}
A linear combination of the Taylor evaluation kernels is
$m^{-1/2}q(\log m)$, where
$q(t)=\sum_jp_j(t)e^{-i\gamma_jt}$. For nonzero coefficients,
independence in Lemma~\ref{v4-arith:mo:exponential-polynomials} gives $q\ne0$.
For $h(x)=x^{-1}|q(\log x)|^2$, the derivative is bounded by
$Cx^{-2}(1+\log x)^d$ for some integer $d$.
This is integrable on $[1,\infty)$. Hence the difference between
$\sum_{m\leq M}h(m)$ and $\int_1^M h(x)\,dx$ stays bounded.
The integral equals $\int_0^{\log M}|q(t)|^2dt$ and tends to
infinity by the positive-half-line estimate in that lemma.
No nonzero combination of the kernels belongs to $\ell^2$.
The orthogonal-complement argument from Theorem~\ref{v4-arith:mo:finite-graph}
now proves dense graph. A linear subspace of a finite-dimensional
space is closed, so density of the range proves surjectivity.

For the last assertion, suppose the kernel were $N$-invariant.
There would be an induced linear operator on the finite-dimensional
range of the residue map. The image of $v_m$ is nonzero: its
last coordinate at any zero is
$b_{\rho,-r_\rho}m^{-\rho}\ne0$.
That image would be an eigenvector with eigenvalue $\log m$.
Infinitely many distinct eigenvalues are impossible in finite
dimension. This contradiction applies without zero-placement
or simplicity assumptions.
\end{proof}

\subsubsection{The Mellin-diagonal dilation generator}

The norm of $\mathscr H_c$ dictates the dilation amplitude. Define
\[
 (V_uf)(x)=e^{cu}f(e^u x),\qquad u\in\mathbb R.
\]
Substitution $y=e^u x$ proves $\|V_uf\|_{\mathscr H_c}=
\|f\|_{\mathscr H_c}$. In logarithmic coordinates, $V_u$ is
translation $g(t)\mapsto g(t+u)$. Its infinitesimal operator on
$\mathscr D$ is $D+c$. With the convention $V_u=e^{-iuP_c}$,
\begin{equation}
 P_c=i(D+c),\qquad
 \mathcal M(P_cf)(s)=i(c-s)F(s).
 \label{v4-arith:mo:dilation}
\end{equation}
In particular $\mathscr H_{1/2}=L^2(dx)$ and
$P_{1/2}=i(x\partial_x+1/2)$. The half-shift comes from the
measure, and the factor $i$ comes from the stated unitary-group
convention. On $L^2(dx/x)$ one has $c=0$ and no half-shift.

\begin{lemma}[the normalized Mellin unitary]
\label{v4-arith:mo:mellin-unitary}
The map
\[
 (\mathcal U_cf)(\tau)=\frac1{\sqrt{2\pi}}
                          \mathcal Mf(c+i\tau)
\]
on $\mathscr D$ extends to a unitary
$\mathscr H_c\to L^2(\mathbb R,d\tau)$. The closure of $P_c$ is
self-adjoint and is multiplication by $\tau$ under this unitary.
Its domain is
\[
 \left\{f\in\mathscr H_c:
       \tau\mathcal U_cf(\tau)\in L^2(d\tau)\right\}.
\]
Equivalently, $U_cf$ has weak derivative in $L^2(dt)$.
\end{lemma}
\begin{proof}
For $g=U_cf$, the displayed map is the Fourier transform
$\widehat g(\tau)=(2\pi)^{-1/2}\int g(t)e^{i\tau t}\,dt$.
Its isometry can be proved directly. For Schwartz functions
$g,h$ and $\epsilon>0$, Fubini and the Gaussian integral give
\[
 \int\widehat g(\tau)\overline{\widehat h(\tau)}
                      e^{-\epsilon\tau^2}\,d\tau
 =\iint g(t)\overline{h(u)}
 \frac{e^{-(t-u)^2/(4\epsilon)}}{2\sqrt{\pi\epsilon}}\,dt\,du.
\]
The Gaussian kernels have integral one and converge as an
approximate identity. Splitting their integrals into
$|t-u|<\delta$ and its complement, then using translation
continuity in $L^2$, proves convergence to $\langle g,h\rangle$.
For $h=g$, monotone convergence proves equality of the two squared
norms. Polarization gives the inner-product identity.
Integration by parts shows that Fourier transforms of Schwartz
functions are Schwartz. The same Gaussian calculation gives
Fourier inversion by first inserting $e^{-\epsilon\tau^2}$ and
then letting $\epsilon$ tend to zero. Thus the range contains the
Schwartz space. Its density and the isometry give a unitary on
$L^2$. Cutoffs and convolution by smooth approximate identities
prove the required density of compact smooth functions.

On compact smooth functions, $\widehat{ig'}=\tau\widehat g$.
Multiplication by $\tau$ on its maximal domain is self-adjoint:
testing against compactly supported functions identifies its
adjoint domain, and division by $\tau\pm i$ gives its resolvents.
The inverse Fourier transform identifies this domain with weakly
differentiable $L^2$ functions whose derivative is in $L^2$.
Such functions can be approximated in the norm
$\|g\|_2+\|g'\|_2$ by cutoffs followed by smooth convolution.
Indeed the cutoff derivative term is bounded by
$C\|g\|_2/L$ for a cutoff varying at scale $L$.
Therefore compact smooth functions form a graph core. This proves
the closure and domain assertions.
\end{proof}

The spectrum of $P_c$ is all of $\mathbb R$ and has no eigenvectors.
A nonzero $L^2$ function cannot be supported at one frequency.
Functions supported in shrinking intervals around any prescribed
frequency give approximate eigenvectors. In contrast, the
occupation operator has the eigenbasis $v_m$ and eigenvalues
$\log m$.

In logarithmic coordinates, $Q=-t$ and $P_c=i\partial_t$.
On compact smooth functions,
\[
 [P_c,Q]=-i\,\mathrm{id}.
\]
With $\hbar=1$, these are a position observable with reversed
orientation and its conjugate momentum. The distribution $Tv_m$
is localized at $t=-\log m$. Its Mellin transform is a plane
wave multiplied by $m^{-c}$ on the line $\operatorname{Re}s=c$;
it is not a frequency point mass. This distinction accounts for
both actions in \eqref{v4-arith:mo:two-actions}.

\subsubsection{Dilation on full residue jets}

Let $\mathscr J_Z=\bigoplus_{\rho\in Z}\mathbb C^{r_\rho}$,
with coordinates $a_{\rho,k}$ as in \eqref{v4-arith:mo:jets}. Set
$a_{\rho,r_\rho}=0$ and define
\begin{equation}
 (K_ca)_{\rho,k}=z_{\rho,c}a_{\rho,k}-i a_{\rho,k+1},
 \qquad z_{\rho,c}=i(c-\rho).
 \label{v4-arith:mo:jet-generator}
\end{equation}
Write $\mathcal R_Zf=(R_{\rho,k}\mathcal Mf)_{\rho,k}$.

\begin{proposition}[the full jet comparison]
\label{v4-arith:mo:jet-comparison}
On $\mathscr D$,
\[
 \mathcal R_ZP_c=K_c\mathcal R_Z.
\]
At a simple zero $\rho=\sigma+i\gamma$, the multiplier is
$z_{\rho,c}=\gamma+i(c-\sigma)$. It equals the ordinate
$\gamma$ exactly when $\sigma=c$. At a zero of multiplicity
$r$, the full block is a Jordan block with eigenvalue $z_{\rho,c}$
and superdiagonal $-i$.
\end{proposition}
\begin{proof}
In \eqref{v4-arith:mo:dilation}, write
$i(c-s)=i(c-\rho)-i(s-\rho)$ and take the residue after
multiplication by $(s-\rho)^k$. The second term is the next
coordinate in \eqref{v4-arith:mo:jets}; it vanishes at the last coordinate.
The superdiagonal is nonzero, so the nilpotent part has exact
order $r$.
\end{proof}

For a double zero, the scalar residue alone obeys
\[
 R_{\rho,0}\mathcal M(P_cf)=
 z_{\rho,c}R_{\rho,0}F-iR_{\rho,1}F.
\]
Dropping the second coordinate drops an actual term. Repeating
$F(\rho)$ in two coordinates instead gives a diagonal action, but
its image has dimension one. It does not retain the full
multiplicity-two jet. Nor does the jet comparison repair the
occupation map: its last coordinate on $Tv_m$ is
$b_{\rho,-r_\rho}m^{-\rho}\ne0$. Intertwining $N$ with $K_c$
would force $\log m=z_{\rho,c}$ for every $m$.

\subsection{Evaluation norms and operator closure}

A residue value is a linear functional on compact smooth tests.
Completing those tests in the squared sum of their residue values
changes the Hilbert space. The change can be measured exactly for
finitely many zeros, without any estimate on the global zero set.

\subsubsection{Finite evaluation and its graph}

\begin{lemma}[independence of exponential polynomials]
\label{v4-arith:mo:exponential-polynomials}
Let $\lambda_1,\ldots,\lambda_n$ be distinct complex numbers.
If polynomials $p_j$ satisfy
$\sum_j p_j(t)e^{\lambda_jt}=0$ on an interval, then all $p_j$ vanish.
A nonzero sum of this form does not belong to $L^2(\mathbb R,dt)$.
\end{lemma}
\begin{proof}
For independence, fix $j$ and apply
$\prod_{k\ne j}(\partial_t-\lambda_k)^{d_k+1}$, where
$d_k\geq\deg p_k$. All terms except the $j$th vanish. On its
polynomial factor this operator is
$\prod_{k\ne j}(\partial_t+\lambda_j-\lambda_k)^{d_k+1}$.
Each factor is invertible on a space of polynomials of bounded
degree: in the monomial basis it is triangular with nonzero
constant diagonal. Therefore $p_j=0$.

For the $L^2$ assertion, choose the largest real part $a$ among
nonzero terms, and the largest polynomial degree $d$ among terms
with that real part. On intervals $[T,T+L]$, divide the sum by
$e^{aT}T^d$. As $T\to+\infty$, its leading part consists of a fixed
finite family $e^{(a+ib_j)u}$, $0\leq u\leq L$, with coefficients
of fixed nonzero magnitudes and phases $e^{ib_jT}$. The Gram
matrix of this family on $[0,L]$ is positive definite by the first
part. All lower-degree or lower-real-part terms tend to zero
uniformly after division. Thus, for sufficiently large $T$,
\[
 \int_T^{T+L}\left|\sum_jp_j(t)e^{\lambda_jt}\right|^2dt
 \geq C e^{2aT}T^{2d}
\]
for some $C>0$. If $a\geq0$, summing over disjoint such intervals
proves divergence. If $a<0$, replace $t$ by $-t$; the largest
real part then becomes positive. This proves the claim.
\end{proof}

\begin{theorem}[finite residue graph]
\label{v4-arith:mo:finite-graph}
For every real $c$ and every nonempty finite zero set $Z$, the map
$\mathcal R_Z:\mathscr D\to\mathscr J_Z$ is onto. Its graph is
dense in $\mathscr H_c\oplus\mathscr J_Z$, where $\mathscr J_Z$
has its Euclidean norm. Hence it is unbounded and not closable
as an operator from $\mathscr H_c$. Its kernel is dense in
$\mathscr H_c$.
\end{theorem}
\begin{proof}
Write $g=U_cf$. The Taylor functionals before the invertible
transformation \eqref{v4-arith:mo:jet-formula} are
\[
 F^{(k)}(\rho)=\int_{\mathbb R}t^k
                       e^{(\rho-c)t}g(t)\,dt.
\]
The kernels are linearly independent by
Lemma~\ref{v4-arith:mo:exponential-polynomials}. No nonzero functional on the
finite-dimensional target can annihilate the range. Thus the
range is the whole target.

Suppose $(u,b)\in L^2(\mathbb R)\oplus\mathscr J_Z$ is orthogonal
to the graph in logarithmic coordinates. For every compact smooth
$g$ this says
\[
 \int g(t)\overline{u(t)}\,dt+\int g(t)q_b(t)\,dt=0,
\]
where $q_b$ is a fixed linear combination of the exponential
polynomials above. Thus $q_b=-\overline u$ as a distribution,
so $q_b\in L^2$. The lemma forces $q_b=0$. Independence and the
invertible jet transformation give $b=0$, and then $u=0$.
A subspace with zero orthogonal complement is dense. The graph
closure is therefore the entire product, which contains
$\{0\}\oplus\mathscr J_Z$ and cannot be a graph.

Choose test functions lifting a basis of $\mathscr J_Z$; their
linear extension is a right inverse $G:\mathscr J_Z\to\mathscr D$.
As a map into $\mathscr H_c$, $G$ is bounded because its domain is
finite-dimensional. Graph density gives, for each $f\in\mathscr H_c$,
$g_n\to f$ and $\mathcal R_Zg_n\to0$. Then
$g_n-G\mathcal R_Zg_n$ lies in the kernel and converges to $f$.
\end{proof}

Define a positive semidefinite form on $\mathscr D$ by
\[
 q_Z(f,h)=\sum_{\rho\in Z}\sum_{k=0}^{r_\rho-1}
 R_{\rho,k}\mathcal Mf\;
 \overline{R_{\rho,k}\mathcal Mh}.
\]
Its nullspace is exactly $\ker\mathcal R_Z$. On the quotient
$\mathscr D/\ker\mathcal R_Z$, this is an inner product.
Surjectivity in Theorem~\ref{v4-arith:mo:finite-graph} identifies its completion
isometrically with $\mathscr J_Z$. The nullspace has infinite
dimension, since $\mathscr D$ has infinite dimension and the
quotient has finite dimension. Thus this positive norm does not
give an injective realization of the original function space.

The form is not closable in $\mathscr H_c$. For a positive form,
closability requires $q(f_n)\to0$ whenever $f_n\to0$ in the
original Hilbert norm and $q(f_n-f_m)\to0$. Dense graph gives
$f_n\to0$ with $\mathcal R_Zf_n$ tending to a fixed nonzero
vector. This sequence violates the form criterion. The same
argument applies to the countable form below whenever its stated
summability hypothesis holds.

\begin{theorem}[symmetry in the finite evaluation norm]
\label{v4-arith:mo:finite-symmetry}
The dilation operator descends to the quotient with form $q_Z$ and
extends to the matrix $K_c$ of \eqref{v4-arith:mo:jet-generator}.
It is symmetric for the Euclidean jet norm exactly when all zeros
in $Z$ are simple and satisfy $\operatorname{Re}\rho=c$.
If either condition fails, no positive definite inner product on
the full jet target can make $K_c$ symmetric.
\end{theorem}
\begin{proof}
The identity $\mathcal R_ZP_c=K_c\mathcal R_Z$ makes the kernel
invariant and proves the induced action. In finite dimension the
induced operator is everywhere defined, bounded, and closed.
If every block has size one and real eigenvalue, its matrix is
real diagonal and hence symmetric.

Conversely, symmetry on a positive inner-product space forces
an eigenvalue $z$ to be real, since
$\langle Kv,v\rangle=z\|v\|^2$ must be real for an eigenvector.
Thus $\operatorname{Re}\rho=c$. A nontrivial Jordan block then
has vectors $u,v\ne0$ with $(K-z)u=v$ and $(K-z)v=0$.
Symmetry would give
$\|v\|^2=\langle(K-z)u,v\rangle=
\langle u,(K-z)v\rangle=0$, a contradiction.
\end{proof}

For one double zero on the line, the matrix is
\[
 \begin{pmatrix}\gamma&-i\\0&\gamma\end{pmatrix}.
\]
Its two real eigenvalues do not remove its nilpotent part. A
positive scalar-residue norm using only $R_{\rho,0}$ generally
has a different defect: its kernel need not be invariant under
$P_c$, because the omitted coordinate $R_{\rho,1}$ appears in
its derivative. The full jets are required before asking for
an induced dilation operator.

\subsubsection{Countable simple zeros and weighted evaluations}

Let $\rho_j=c+i\gamma_j$ be a countable set of distinct simple
zeros, with no finite accumulation. Choose positive weights $w_j$.
The raw residue norm would be
\begin{equation}
 q_w(f)=\sum_j w_j
       \left|\frac{F(\rho_j)}{\Xi'(\rho_j)}\right|^2.
 \label{v4-arith:mo:weighted-residue-norm}
\end{equation}
Put $\mu_j=w_j/|\Xi'(\rho_j)|^2$. Multiplication of the residue
coordinates by $\Xi'(\rho_j)$ identifies the target norm with
$\ell^2(\mu)$, whose norm is $\sum_j\mu_j|a_j|^2$.
This conversion retains the residue normalization, including its
complex phase.

A sufficient summability condition is
\begin{equation}
 \sum_j\mu_j(1+\gamma_j^2)^{-b}<\infty
 \quad\text{for some }b\geq0.
 \label{v4-arith:mo:weight-condition}
\end{equation}
It is a hypothesis on the chosen weights and the actual divisor.
It is not inferred for the raw choice $w_j=1$. Under
\eqref{v4-arith:mo:weight-condition}, define
\[
 E:\mathscr D\longrightarrow\ell^2(\mu),\qquad
 (Ef)_j=F(c+i\gamma_j),\qquad
 (Ka)_j=\gamma_ja_j.
\]
The maximal domain of $K$ consists of sequences with
$\sum_j\mu_j\gamma_j^2|a_j|^2<\infty$. Compact smooth tests have
Schwartz Mellin transforms on the vertical line, by repeated
integration by parts in $t$. Therefore $Ef\in\operatorname{Dom}K$
and $EP_c=KE$ on $\mathscr D$.

\begin{theorem}[the atomic completion and the original norm]
\label{v4-arith:mo:countable-completion}
Under \eqref{v4-arith:mo:weight-condition}, $E\mathscr D$ is dense in
$\ell^2(\mu)$ and is a graph core for $K$.
The completion of $\mathscr D/\ker E$ in the norm $\|Ef\|$
is therefore $\ell^2(\mu)$. The induced dilation has self-adjoint
closure $K$ on its stated maximal domain.

As a map from the original space $\mathscr H_c$, $E$ has dense
graph in $\mathscr H_c\oplus\ell^2(\mu)$. In particular it is
not closable and has no closable extension agreeing on
$\mathscr D$.
\end{theorem}
\begin{proof}
Fix $j_0$ and use $f_L$ from \eqref{v4-arith:mo:continuous-sequence} with
$\gamma_0=\gamma_{j_0}$. Put
$\Psi(y)=\int\psi(t)e^{iyt}\,dt$. Then
\[
 (Ef_L)_j=\Psi\bigl(L(\gamma_j-\gamma_{j_0})\bigr).
\]
The central value is one and every other value tends to zero.
For $L\geq1$, rapid decay of $\Psi$ gives a uniform bound
$C_n(1+|\gamma_j-\gamma_{j_0}|)^{-n}$ for every $n$.
Choosing $n$ sufficiently large in
\eqref{v4-arith:mo:weight-condition} gives a summable majorant even after
multiplication by $1+\gamma_j^2$. Dominated convergence proves
\[
 Ef_L\longrightarrow e_{j_0},\qquad
 KEf_L\longrightarrow\gamma_{j_0}e_{j_0}
 \quad\text{in }\ell^2(\mu),
\]
while $f_L\to0$ in $\mathscr H_c$.
Finite-coordinate vectors form a graph core for the maximal
real diagonal operator $K$, by truncation of the two convergent
weighted sums. The approximations just constructed show that
$E\mathscr D$ is also a graph core. The diagonal adjoint-domain
calculation proves self-adjointness, as in
Lemma~\ref{v4-arith:mo:mellin-unitary}.

The graph closure of $E$ contains $(0,e_j)$ for every $j$, hence
$\{0\}\oplus\ell^2(\mu)$. It also contains $(f,Ef)$ for every
test $f$. Subtracting the vertical vectors and using density of
$\mathscr D$ in $\mathscr H_c$ gives the whole product. This
proves the last assertion.
\end{proof}

The hypotheses of this theorem can always be imposed on a selected
countable critical-line simple divisor by changing the weights.
For example, fix an enumeration and set
\[
 \mu_j=2^{-j}(1+\gamma_j^2)^{-j},\qquad
 w_j=|\Xi'(\rho_j)|^2\mu_j.
\]
All polynomial moments of $\mu$ are finite: for a fixed degree,
only finitely many initial terms need separate consideration.
This produces a positive evaluation completion. The weights
are chosen data; they are not the raw residue weights or a native
arithmetic pairing.

For $w_j=1$, the exact domain of \eqref{v4-arith:mo:weighted-residue-norm}
is
\[
 \mathscr D_{\mathrm{raw}}=
 \left\{f\in\mathscr D:
 \sum_j|F(\rho_j)/\Xi'(\rho_j)|^2<\infty\right\}.
\]
No assertion that this equals $\mathscr D$, is dense in
$\mathscr H_c$, or satisfies \eqref{v4-arith:mo:weight-condition} is made
without bounds on the inverse derivatives. Even existence of an
individual scalar residue does not imply convergence of this sum.
If \eqref{v4-arith:mo:weight-condition} is established for the raw weights,
Theorem~\ref{v4-arith:mo:countable-completion} applies and still gives
nonclosability in the original norm. Without that estimate, a
full vector-valued closure claim remains separate from the
finite-coordinate obstruction.

\begin{proposition}[bounded spectral transfer]
\label{v4-arith:mo:bounded-transfer}
A bounded map intertwining the unitary groups of $P_c$ on
$\mathscr H_c$ and $K$ on $\ell^2(\mu)$ is zero, in either
direction.
\end{proposition}
\begin{proof}
In the direction from the atomic space, the image of $e_j$ would
be an eigenvector of $P_c$ with eigenvalue $\gamma_j$.
Under $\mathcal U_c$, group covariance for all rational times
forces its support to lie in $\{\gamma_j\}$, a null set.
Thus it is zero, and density of the finite-coordinate vectors
finishes the argument. Taking the adjoint gives the opposite
direction, since both groups are unitary.
\end{proof}

\begin{proposition}[occupation symmetry in a sampling norm]
\label{v4-arith:mo:occupation-positive-obstruction}
Let $\rho_j=c+i\gamma_j$ be a nonempty countable set of simple
zeros, with positive weights $\mu_j$ and
$\sum_j\mu_j<\infty$. The evaluation form
\[
 q(a,b)=\sum_j\mu_jF_a(\rho_j)\overline{F_b(\rho_j)}
 \quad(a,b\in\mathscr E)
\]
is well-defined and positive semidefinite. The occupation operator
$N$ is not symmetric for this form on its finite-support core.
\end{proposition}
\begin{proof}
Every finite Dirichlet polynomial is bounded on a fixed vertical
line. The summability assumption therefore defines the form.
Put $\phi(t)=\sum_j\mu_je^{i\gamma_jt}$. Uniform convergence
makes $\phi$ continuous, and $\phi(0)=\sum_j\mu_j>0$.
The coordinate pairing is
\[
 q(v_m,v_n)=(mn)^{-c}\phi(\log(n/m)).
\]
Symmetry would require
$(\log m-\log n)q(v_m,v_n)=0$ for all $m,n$.
Taking $n=m+1$ would give
$\phi(\log(1+1/m))=0$. Continuity at zero contradicts
$\phi(0)>0$.
\end{proof}

Thus even a positive critical-line evaluation completion cannot
make occupation energy symmetric on this core. The same norm
does make the induced continuous dilation self-adjoint under
Theorem~\ref{v4-arith:mo:countable-completion}. These statements concern different
source operators, with the exact actions in \eqref{v4-arith:mo:two-actions}.

\subsubsection{A countable completion with all principal parts}

The full principal parts also admit a Hilbert completion when
multiplicities are unbounded or zeros lie off the chosen line.
The construction below chooses the norm after the divisor and
finite interpolation functions have been specified.

Enumerate distinct zeros as $\rho_j$, $j\geq1$, with
$|\rho_j|\to\infty$, and let $r_j$ be their finite multiplicities.
Fix real $c$. Write $R_jf\in\mathbb C^{r_j}$ for the full residue
vector of $\mathcal Mf$ and set
\[
 z_j=i(c-\rho_j),\qquad K_j=z_jI-iS_j,
 \qquad (S_ja)_k=a_{k+1},\quad a_{r_j}=0.
\]
Each block has its Euclidean norm, so $\|S_j\|\leq1$.
For positive numbers $w_j$ define
\[
 \mathscr J_w=\left\{a=(a_j):
              \sum_jw_j\|a_j\|^2<\infty\right\},\qquad
 (K_wa)_j=K_ja_j.
\]
Its maximal domain is
\begin{equation}
 \operatorname{Dom}K_w=
 \left\{a\in\mathscr J_w:
             \sum_jw_j|z_j|^2\|a_j\|^2<\infty\right\}.
 \label{v4-arith:mo:full-domain}
\end{equation}
Indeed $K_w$ differs from multiplication by $z_j$ by the bounded
operator $-i\bigoplus_jS_j$. Equivalently its domain requires
$\sum_jw_j\|K_ja_j\|^2<\infty$.

\begin{theorem}[weighted full-jet completion]
\label{v4-arith:mo:all-jets}
For every such zero divisor and every real $c$, positive weights
$w_j$ can be chosen with the following properties. The map
$R:f\mapsto(R_jf)_j$ sends $\mathscr D$ into
$\operatorname{Dom}K_w$ and satisfies $RP_c=K_wR$.
Its range is dense and is a graph core for $K_w$.
The completion of $\mathscr D/\ker R$ in the norm $\|Rf\|$
is $\mathscr J_w$, and the induced dilation closes to $K_w$.
The graph of $R$ in $\mathscr H_c\oplus\mathscr J_w$ is dense.
Thus $R$ and its positive evaluation form are not closable in
the original norm.

The closed operator $K_w$ has compact resolvent and spectrum
$\{z_j:j\geq1\}$. The generalized eigenspace at $z_j$ is one
block of dimension $r_j$. It is self-adjoint exactly when all
$r_j=1$ and $\operatorname{Re}\rho_j=c$.
\end{theorem}
\begin{proof}
Let $I_n=\{(j,k):j\leq n,\ 0\leq k<r_j\}$ and denote its
coordinate basis by $e_h$. The finite graph theorem supplies,
for each $n$ and $h\in I_n$, a test $f_{n,h}$ such that
\begin{equation}
 \|f_{n,h}\|_{\mathscr H_c}\leq n^{-1},\qquad
 (R_jf_{n,h})_{j\leq n}=e_h.
 \label{v4-arith:mo:small-lifts}
\end{equation}
Here exact interpolation follows from graph density by correction
with the bounded finite-dimensional right inverse used in its proof.
All these choices precede the choice of weights.

For $g=U_cf$ supported in $[-m,m]$, each residue is an integral
against an exponential polynomial. Hence there is a finite
constant $C_{j,m}$ with
$\|R_jf\|\leq C_{j,m}\|g\|_2$ on that support.
One explicit choice is
\[
 C_{j,m}^2=\sum_{k=0}^{r_j-1}\int_{-m}^m
 \left|e^{(\rho_j-c)t}
 \sum_{\ell=0}^{r_j-k-1}
 b_{\rho_j,-k-\ell-1}\frac{t^\ell}{\ell!}\right|^2dt.
\]
Cauchy--Schwarz proves the bound. Define the finite maximum
\[
 A_j=\max\left(
 1,\ \max_{1\leq m\leq j}C_{j,m}^2,
 \max_{\substack{1\leq n<j\\h\in I_n}}
                  n^2\|R_jf_{n,h}\|^2\right),
 \qquad
 w_j=\frac{2^{-j}}{(1+|z_j|^2)A_j}.
\]
The last maximum is omitted for $j=1$.
For any fixed compact smooth test and sufficiently large $j$,
the summand $w_j(1+|z_j|^2)\|R_jf\|^2$ is at most
$2^{-j}\|U_cf\|_2^2$. This proves membership in the domain.
The coordinate identity from Proposition~\ref{v4-arith:mo:jet-comparison} therefore
holds as an identity in $\mathscr J_w$.

For a fixed coordinate $h$, take $n$ large enough that $h\in I_n$.
The first $n$ blocks of $Rf_{n,h}-e_h$ vanish. The remaining
blocks satisfy
\[
 \sum_{j>n}w_j(1+|z_j|^2)\|R_jf_{n,h}\|^2
 \leq n^{-2}\sum_{j>n}2^{-j}.
\]
Since $\|K_j\|\leq |z_j|+1$, this gives
$Rf_{n,h}\to e_h$ and $K_wRf_{n,h}\to K_we_h$.
Finite block truncations are a graph core for the maximal operator,
by convergence of the two sums defining its graph norm.
Thus $R\mathscr D$ is a graph core and is dense.
Equation \eqref{v4-arith:mo:small-lifts} also gives $f_{n,h}\to0$.
The graph closure contains $\{0\}\oplus\mathscr J_w$.
Subtracting its vertical vectors from $(f,Rf)$ and using density
of $\mathscr D$ proves density in the whole product.
The quotient and nonclosability assertions follow.

The maximal block operator is closed: convergence in its graph
implies convergence in every finite block, where each matrix is
closed. Its adjoint is the maximal operator with blocks
$\overline z_jI+iS_j^*$, as testing against finite block vectors
shows. These formulas also prove self-adjointness when all blocks
are real scalars. If it is symmetric, its restriction to every
finite block is symmetric. Theorem~\ref{v4-arith:mo:finite-symmetry} forces every
block to be a real scalar.

Choose $\lambda$ different from every $z_j$. For all sufficiently
large $j$, $|z_j-\lambda|>2$. The finite geometric series for
$(K_j-\lambda)^{-1}$ then gives
\[
 \|(K_j-\lambda)^{-1}\|
 \leq\frac1{|z_j-\lambda|-1}\longrightarrow0.
\]
The finitely many other inverses are bounded. Their direct sum
maps $\mathscr J_w$ into \eqref{v4-arith:mo:full-domain}, since
$K_ja_j=y_j+\lambda a_j$ for
$a_j=(K_j-\lambda)^{-1}y_j$.
Thus $\lambda$ belongs to the resolvent. Truncating this inverse
to finitely many finite-dimensional blocks converges in operator
norm, so the resolvent is compact. Each $z_j$ is an eigenvalue.
On every other block $K_l-z_j$ is invertible, so its generalized
eigenspace is precisely the $j$th block.
\end{proof}

The theorem requires no bound on the multiplicities. The norm
depends on the Laurent coefficients and the chosen interpolation
functions. Making a weight smaller preserves every estimate and
every graph approximation in the proof. For a divisor invariant
under $\rho\mapsto2c-\overline\rho$, one may therefore replace
the two weights on each orbit by their minimum. The resulting
residue involution is antiunitary. This does not make a nontrivial
Jordan block symmetric.

For a double zero on the line, the ordinary trace of a polynomial
in $K_j$ is $2p(\gamma_j)$, while the operator retains its
nonzero nilpotent part. More generally,
\[
 \operatorname{tr}p(K_j)=r_jp(z_j).
\]
This follows by triangularity. Consequently these traces cannot
distinguish the full block from $r_j$ repeated scalar eigenvalues.
Equality of spectral traces alone does not construct a map of
their full jet representations.


\subsubsection{Stronger source norms and boundary values}

Evaluation can be continuous after the source topology changes.
On finite occupation sequences define, for real $a$,
\[
 \|u\|_a^2=\sum_{m\geq1}(1+(\log m)^2)^a|u_m|^2,
\]
and let $\mathscr H_{\mathrm{occ},a}$ be the corresponding weighted
$\ell^2$ space. Fix $\rho=1/2+i\gamma$ and an integer $k\geq0$.

\begin{proposition}[the exact logarithmic weight threshold]
\label{v4-arith:mo:weight-threshold}
The functional
$u\mapsto F_u^{(k)}(\rho)=\sum_m u_m(-\log m)^k m^{-\rho}$
on finite sequences is bounded for $\|\cdot\|_a$ exactly when
$a>k+1/2$. If $\operatorname{Re}\rho>1/2$, it is bounded for
every real $a$. If $\operatorname{Re}\rho<1/2$, no real $a$
makes it bounded.
\end{proposition}
\begin{proof}
The squared dual norm of a finite-coordinate functional increases
under truncation to
\[
 \sum_m\frac{(\log m)^{2k}m^{-2\operatorname{Re}\rho}}
                    {(1+(\log m)^2)^a}.
\]
Cauchy--Schwarz proves sufficiency. Truncated coefficient vectors
proportional to the weighted conjugate kernel attain the finite
dual norms, proving necessity. At real part $1/2$, integral
comparison with $t=\log x$ gives the tail integral
$\int^\infty t^{2k-2a}\,dt$. It is finite exactly when
$a>k+1/2$. To the right, exponential decay in $t$ dominates every
power. To the left, exponential growth defeats every such power.
\end{proof}

At a zero of order $r$, the scalar residue functional on finite
occupation sequences has highest derivative order $r-1$, with
nonzero coefficient. On the critical line its exact threshold
is therefore $a>r-1/2$. The same condition controls all full jet
coordinates. Lower derivative terms cannot cancel the nonzero
leading polynomial in $\log m$.

These extensions are boundary-value functionals of the Dirichlet
series from Proposition~\ref{v4-arith:mo:dirichlet-domain}. Their defining sums converge
absolutely under the stated weight condition. They do not supply
holomorphic continuation across the boundary. Calling them actual
meromorphic residues for an infinite sequence additionally requires
such continuation. Nor do stronger norms repair
Proposition~\ref{v4-arith:mo:occupation-defect}, which already fails on finite sequences.

The occupation operator remains self-adjoint in
$\mathscr H_{\mathrm{occ},a}$ on its maximal weighted domain,
because both the norm and the operator are diagonal in $v_m$.
For continuous tests, a different change is needed to bound point
evaluation: put $\|g\|_{a,t}^2=\int(1+t^2)^a|g(t)|^2dt$.
The same dual-integral argument bounds the $k$th Fourier evaluation
exactly when $a>k+1/2$. This norm is not translation-invariant.
For $w(t)=(1+t^2)^a$ and compact smooth $g,h$,
\[
 \langle i g',h\rangle_w-\langle g,i h'\rangle_w
 =-i\int w'(t)g(t)\overline{h(t)}\,dt.
\]
Thus the same dilation generator need not remain symmetric in
this stronger continuous norm. A topology that controls residues
must be checked against the required operator and pairing.

\subsection{Critical-line symmetry and the arithmetic target}

An involution can preserve residue data without making the dilation
operator symmetric. This distinction is visible in the signs of
full principal parts.

\subsubsection{Inversion and conjugation}

For real $c$, define on $\mathscr H_c$
\[
 (J_cf)(x)=x^{-2c}\overline{f(x^{-1})}.
\]
In logarithmic coordinates,
$U_cJ_cf(t)=\overline{U_cf(-t)}$. Consequently $J_c$ is an
antiunitary involution. Direct substitution in the Mellin integral
gives
\begin{equation}
 \mathcal M(J_cf)(s)=\overline{F(2c-\overline s)},
 \qquad J_cP_c=P_cJ_c
 \label{v4-arith:mo:inversion}
\end{equation}
on compact smooth tests, and then on the closed operator domain.
For the domain assertion, logarithmic reflection and conjugation
preserve the weak-derivative norm from
Lemma~\ref{v4-arith:mo:mellin-unitary}.

Suppose now that the entire function satisfies
\begin{equation}
 \Xi(2c-\overline s)=\overline{\Xi(s)}.
 \label{v4-arith:mo:entire-symmetry}
\end{equation}
Write $\rho^\#=2c-\overline\rho$. Its zeros and multiplicities
are preserved by $\rho\mapsto\rho^\#$. Choose the finite zero set
$Z$ invariant under this map.

\begin{proposition}[the residue involution]
\label{v4-arith:mo:residue-involution}
On full jet coordinates, inversion acts by
\[
 (J_Za)_{\rho,k}=(-1)^{k+1}\overline{a_{\rho^\#,k}}.
\]
Thus $\mathcal R_ZJ_c=J_Z\mathcal R_Z$, and $J_Z$ is an
antiunitary involution for the Euclidean jet norm. It commutes
with $K_c$. At a simple critical-line zero, the residue action
is minus complex conjugation.
\end{proposition}
\begin{proof}
Put $s=\rho+z$. Equations \eqref{v4-arith:mo:inversion} and
\eqref{v4-arith:mo:entire-symmetry} identify the quotient for $J_cf$ with
the conjugate of the quotient for $f$ at
$\rho^\#-\overline z$. A coefficient of order $-(k+1)$ therefore
acquires the factor $(-1)^{k+1}$. This proves the formula.
The norm and involution assertions follow by permuting coordinates
and conjugating them. For commutation, note that
$z_{\rho^\#,c}=\overline{z_{\rho,c}}$ and that consecutive signs
in the formula are opposite. Substitution into
\eqref{v4-arith:mo:jet-generator} gives $J_ZK_c=K_cJ_Z$.
\end{proof}

The same assertion holds for positive diagonal jet weights invariant
under $\rho\mapsto\rho^\#$. Commutation with this antiunitary does
not remove a Jordan block. The symmetry criterion of
Theorem~\ref{v4-arith:mo:finite-symmetry} still requires simple zeros on the line.

The involution extends to compactly supported distributions by
\[
 \langle J_cu,\phi\rangle
 =\overline{\left\langle u,
         x^{2c-2}\overline{\phi(x^{-1})}\right\rangle}.
\]
Changing variables proves agreement with its action on functions.
For occupation distributions,
\[
 J_cTv_m=m^{1-2c}\delta_m.
\]
At $c=1/2$ this is $\delta_m$, with Mellin transform $m^{s-1}$.
Its $Q$ eigenvalue is $-\log m$. Except for $m=1$, it lies outside
the distributional image supported at the reciprocal integers.
Thus an inversion-invariant space must contain both signs of this
logarithmic position observable. Identifying this involution with
an arithmetic modular conjugation requires a separate map of the
actual spaces and pairings.

\subsubsection{The completed Riemann function}

The preceding results apply to any nonzero entire $\Xi$.
For the arithmetic specialization, an integral formula fixes the
normalization and both symmetries without a choice of zero locations.
For $t>0$, put
\[
 \vartheta(t)=\sum_{n\in\mathbb Z}e^{-\pi n^2t},\qquad
 \psi(t)=\sum_{n\geq1}e^{-\pi n^2t}.
\]
These series and their derivatives converge uniformly on compact
subsets of $t>0$.

The Gaussian identity
$\int_{\mathbb R}e^{-\pi t x^2}e^{-2\pi i yx}\,dx
=t^{-1/2}e^{-\pi y^2/t}$ follows by differentiation in $y$ and
integration by parts; its value at zero is the Gaussian integral.
Periodize $e^{-\pi t x^2}$ over the integers. Its Fourier coefficients
are the displayed Gaussian values at integer $y$. Both the
periodized series and the resulting Fourier series converge
absolutely and uniformly. Equality follows from uniqueness of
Fourier coefficients for continuous periodic functions. For
completeness, uniqueness follows by convolution with
\[
 F_n(x)=\frac1n\left|\sum_{k=0}^{n-1}e^{2\pi i kx}\right|^2:
\]
these nonnegative kernels have integral one and their mass outside
any neighborhood of the integers tends to zero. Their convolutions
converge uniformly to a continuous periodic function. Evaluating
at zero gives
\[
 \vartheta(t)=t^{-1/2}\vartheta(t^{-1}),\qquad
 \psi(t)=t^{-1/2}\psi(t^{-1})+\tfrac12(t^{-1/2}-1).
\]

Define
\begin{equation}
 \xi_{\mathbb R}(s)=\frac12+
 \frac{s(s-1)}2\int_1^\infty
 \psi(t)\bigl(t^{s/2}+t^{(1-s)/2}\bigr)\,\frac{dt}{t}.
 \label{v4-arith:mo:xi-integral}
\end{equation}
The exponential decrease of $\psi$ makes this integral entire,
with every derivative locally uniformly convergent. It gives
$\xi_{\mathbb R}(0)=\xi_{\mathbb R}(1)=1/2$ and
\[
 \xi_{\mathbb R}(1-s)=\xi_{\mathbb R}(s),\qquad
 \xi_{\mathbb R}(\overline s)=\overline{\xi_{\mathbb R}(s)}.
\]
In particular it satisfies \eqref{v4-arith:mo:entire-symmetry} at $c=1/2$.

On $\operatorname{Re}s>1$, define
$\zeta(s)=\sum_{n\geq1}n^{-s}$ and
$\Gamma(s/2)=\int_0^\infty e^{-u}u^{s/2-1}\,du$.
Absolute convergence permits summation under the integral and gives
\[
 \int_0^\infty\psi(t)t^{s/2}\,\frac{dt}{t}
       =\pi^{-s/2}\Gamma(s/2)\zeta(s).
\]
Split the integral at one and use the theta identity on $(0,1)$.
The elementary contribution is
$1/(s-1)-1/s=1/(s(s-1))$; the remaining integral is that in
\eqref{v4-arith:mo:xi-integral}. Therefore
\[
 \xi_{\mathbb R}(s)=
 \tfrac12s(s-1)\pi^{-s/2}\Gamma(s/2)\zeta(s)
 \quad(\operatorname{Re}s>1).
\]
Equation \eqref{v4-arith:mo:xi-integral} supplies the entire continuation.
No assertion about simplicity or critical-line placement of its
zeros was used.

\begin{proposition}[the occupation vacuum and the zero divisor]
\label{v4-arith:mo:vacuum}
One has $1/4<\xi_{\mathbb R}(1/2)<1/2$.
Suppose a self-adjoint operator $H$ has the property that every
eigenvalue of $1/2+iH$ is a zero of $\xi_{\mathbb R}$.
There is no injective linear map $A$ whose domain contains $v_1$,
whose image of $v_1$ belongs to $\operatorname{Dom}H$, and which
satisfies $ANv_1=HAv_1$.
\end{proposition}
\begin{proof}
At $s=1/2$, \eqref{v4-arith:mo:xi-integral} reads
\[
 \xi_{\mathbb R}(1/2)=\frac12-
          \frac14\int_1^\infty\psi(t)t^{-3/4}\,dt.
\]
For $t\geq1$, the inequality $n^2\geq n$ gives
$0<\psi(t)\leq e^{-\pi t}/(1-e^{-\pi})$.
Consequently the positive integral is bounded above by
$e^{-\pi}/[\pi(1-e^{-\pi})]<1$. This proves the bounds.
Since $Nv_1=0$, the proposed identity gives $HAv_1=0$.
Injectivity makes $Av_1\ne0$, so $1/2$ would be an eigenvalue
of $1/2+iH$ and a zero of $\xi_{\mathbb R}$. The bounds exclude it.
\end{proof}

For $\Xi=\xi_{\mathbb R}$ and $c=1/2$, the full jet eigenvalue is
\[
 z_{\rho,1/2}=-i(\rho-1/2)
       =\operatorname{Im}\rho+i(1/2-\operatorname{Re}\rho).
\]
The affine shift $1/2+iK_{1/2}$ has eigenvalues $\rho$ and, on
full jets, superdiagonal $1$. Replacing $z_{\rho,1/2}$ by
$\operatorname{Im}\rho$ projects the eigenvalue to the critical
line. Removing the superdiagonal removes the generalized-eigenvector
data at multiple zeros. Both changes require an explicit comparison.

If the selected zeros are simple and on the critical line, their
weighted evaluation completion from
Theorem~\ref{v4-arith:mo:countable-completion} has a self-adjoint diagonal dilation
operator. Its spectrum and norm were specified through the selected
zero data and weights. This conditional construction does not
produce an independently defined arithmetic positive pairing or
establish critical-line placement. With full multiple-zero jets,
Theorem~\ref{v4-arith:mo:finite-symmetry} gives an additional obstruction to a
positive symmetric realization.

\subsubsection{Local quotient representations and traces}

Let $\mathcal O_\rho$ be the algebra of holomorphic germs at
$\rho$. If $\Xi$ has a zero of order $r$ there, the linear map
\[
 \mathcal O_\rho/(\Xi)\longrightarrow\mathbb C^r,
 \qquad [F]\longmapsto(R_{\rho,k}F)_{0\leq k<r}
\]
is an isomorphism. Indeed $F/\Xi$ has no principal part exactly
when $F$ is divisible by $\Xi$ as a germ. Taylor reduction modulo
$(s-\rho)^r$ and \eqref{v4-arith:mo:jet-formula} prove surjectivity.
Multiplication by $s$ becomes $\rho I+S$, where
$(Sa)_k=a_{k+1}$. Thus the affine operator $c+iK_c$ is exactly
the local multiplication operator. This comparison preserves
multiplicity and its nilpotent action.

This local quotient is the algebraic form of zero multiplicity.
Meyer constructs a global virtual representation of the idele
class group on nuclear bornological spaces. Its character is
the Weil distribution, and its spectral multiplicities are the
pole orders of the completed $L$-function
(Meyer, arXiv:math/0311468v3, Theorems~5.8 and~5.11, pp.~43 and~47).
Zeros contribute through the negative representation. That result
does not specify the positive evaluation weights constructed here.
An identification with these Hilbert spaces requires maps from
the global representation, continuity in both norms, and compatibility
with the scaling action and the chosen pairing.

For finite blocks the distinction between trace and representation
is explicit. The trace $\operatorname{tr}p(\rho I+S)=rp(\rho)$
equals the trace on $r$ scalar copies of $\rho$.
When $r>1$ these representations are not isomorphic: the first has
a nonzero nilpotent part, while the second has none.
For infinite blocks an ordinary operator trace additionally requires
trace-class convergence. In the simple critical-line case,
\[
 \operatorname{Tr}\varphi(K_w)=\sum_j\varphi(\gamma_j)
 \quad\text{if}\quad \sum_j|\varphi(\gamma_j)|<\infty.
\]
To prove this formula, use the orthonormal basis
$w_j^{-1/2}e_j$ and the diagonal trace-class criterion.
In particular the weights disappear from the trace.
Hence this trace cannot recover the positive form chosen on
the test functions.

The occupation operator has a different trace. For real $\beta>1$,
\[
 \operatorname{Tr}e^{-\beta N}=\sum_{m\geq1}m^{-\beta}
 =\zeta(\beta).
\]
The positive diagonal sum converges by integral comparison and
diverges for $\beta\leq1$. A trace over zero ordinates involves
a different index set and requires a proved comparison of the
operators and trace functionals. Mellin evaluation and equality
of local multiplicities do not supply that comparison.

\subsubsection{The occupation realization that remains}

An occupation-to-zero realization would require a Hilbert space
$\mathscr H_\xi$, an independently specified positive inner product,
and a self-adjoint operator $H$ there. It would also require an
injective linear map $A$ on a dense $N$-invariant domain
$\mathscr D_N\subset\operatorname{Dom}N$, with dense range in
$\mathscr H_\xi$, such that
\[
 A\mathscr D_N\subset\operatorname{Dom}H,\qquad
 AN=HA\quad\text{on }\mathscr D_N.
\]
Closedness and domain compatibility are part of this operator
problem. Matching the arithmetic zero divisor further requires
the eigenvalues of $1/2+iH$, with their multiplicities, to give
the actual zeros of $\xi_{\mathbb R}$. Any determinant or trace
comparison requires its own definition, convergence domain, and
normalization.

The explicit map $T$ constructs the occupation-to-distribution
edge on every finite sequence. Its residue composition fails the
required energy identity by Proposition~\ref{v4-arith:mo:occupation-defect}, including
multiple zeros. A domain containing the finite occupation core
cannot fix that algebraic failure by completion. Domains excluding
that core are different operator problems and require new maps.
Proposition~\ref{v4-arith:mo:vacuum} also excludes every injective map
with the actual zero-divisor target when its domain contains $v_1$.
The normalized dilation operator instead has the exact jet
comparison Proposition~\ref{v4-arith:mo:jet-comparison}; it acts on a continuous
spectral carrier. Its evaluation map is not a closable transfer
from the original Hilbert space in the settings proved above.

The physical distinction is between logarithmic position and its
conjugate momentum. Integer occupation selects discrete position
values. Mellin dilation measures continuous momentum. Residue
evaluation samples a function of that momentum at a chosen divisor;
a positive sampling norm changes the state space. The displayed
maps and obstructions specify these operations. They supply no
identification of their thermal traces or of an arithmetic trace
with the sampled norm.

